% glossary.tex — \input dari report.tex (lampiran)
\chapter{Daftar Istilah}
\label{ch:glosarium}

Lampiran ini menjelaskan istilah-istilah teknis yang dipakai di laporan, disertai
contoh singkat. Istilah teknis sengaja dipertahankan dalam bahasa aslinya karena itulah
bentuk yang lazim dipakai praktisi. Daftar ini panjang karena topiknya menjembatani dua
ekosistem (BEAM/Elixir dan Kubernetes)---namun \textbf{pembaca tidak perlu menghafal
semuanya}: cukup pahami dua belas istilah inti pada Tabel~\ref{tab:inti}, sisanya bersifat
rujukan saat diperlukan.

\section*{Peta Cepat: Istilah Inti}
\begin{table}[H]\centering\small
\caption{Dua belas istilah inti yang cukup untuk mengikuti alur laporan.}
\label{tab:inti}
\begin{tabular}{@{}p{0.30\linewidth}p{0.62\linewidth}@{}}
\toprule
\textbf{Istilah} & \textbf{Inti} \\
\midrule
BEAM & VM Erlang tempat Elixir berjalan. \\
OTP / \emph{supervision tree} & Pola supervisi yang menyalakan ulang \emph{process} yang gagal. \\
libcluster & Pembentukan \emph{cluster} BEAM otomatis di Kubernetes. \\
Horde & \emph{Supervisor}/\emph{registry} terdistribusi dengan \emph{handoff}. \\
Handoff & Pemindahan \emph{process} beserta \emph{state} ke node lain. \\
Split-brain & Partisi jaringan yang membuat dua sisi \emph{cluster} sama-sama aktif. \\
Probe (\emph{liveness}/\emph{readiness}) & Pemeriksaan kesehatan oleh Kubernetes. \\
Grace period & Tenggat antara \texttt{SIGTERM} dan \texttt{SIGKILL}. \\
PDB & Batas \emph{pod} yang boleh \emph{down} saat gangguan sukarela. \\
Operator (Bonny) & Pengendali kustom Kubernetes, ditulis dalam Elixir. \\
Kubernetes Lease & Objek koordinasi/\emph{fencing} berbasis etcd. \\
Kuorum BEAM & Jumlah node minimum agar \emph{cluster} otoritatif. \\
\bottomrule
\end{tabular}
\end{table}

\section{Ekosistem BEAM / Elixir}
\begin{description}[style=nextline,leftmargin=1.4em]
  \item[BEAM] Mesin virtual Erlang tempat Elixir dijalankan; menjalankan jutaan
    \emph{process} ringan terisolasi. \emph{Contoh:} aplikasi Phoenix berjalan sebagai
    satu instans BEAM dalam satu \emph{pod}.
  \item[OTP] Pustaka dan pola Erlang (\emph{supervisor}, \texttt{GenServer}) untuk sistem
    tahan-gagal. \emph{Contoh:} \emph{supervisor} OTP menyalakan ulang pekerja yang \emph{crash}.
  \item[Supervision tree] Hierarki \emph{supervisor}--anak yang menentukan pola pemulihan.
    \emph{Contoh:} lihat Listing~\ref{lst:sup}.
  \item[Supervisor] \emph{Process} yang memantau dan menyalakan ulang anak menurut
    strategi (\texttt{:one\_for\_one}, dll). \emph{Contoh:} memulai ulang hanya anak yang gagal.
  \item[GenServer] Abstraksi \emph{process} \emph{stateful}. \emph{Contoh:} satu sesi
    dokumen direpresentasikan sebagai satu \texttt{GenServer}.
  \item[Process (BEAM)] Unit eksekusi ringan BEAM, bukan \emph{process} OS.
    \emph{Contoh:} ratusan ribu \emph{process} dalam satu VM.
  \item[Let it crash] Filosofi membiarkan \emph{process} gagal lalu dipulihkan supervisor.
    \emph{Contoh:} tidak menulis \emph{defensive code}, cukup biarkan \emph{crash} dan restart.
  \item[Restart] Penyalaan ulang \emph{process} oleh supervisor setelah \emph{crash}.
    \emph{Contoh:} \emph{decoder} yang \emph{crash} disalakan ulang otomatis.
  \item[Restart-budget] Batas intensitas \emph{restart} (\texttt{max\_restarts} dalam
    \texttt{max\_seconds}); bila terlampaui, supervisor menyerah dan tereskalasi ke induk.
    \emph{Contoh:} \texttt{3} \emph{restart} dalam \texttt{5} detik.
  \item[libcluster] Pembentukan \emph{cluster} BEAM otomatis di Kubernetes. \emph{Contoh:}
    strategi \texttt{Kubernetes.DNS} menemukan \emph{peer} lewat \emph{headless service}.
  \item[Horde] \emph{Supervisor}/\emph{registry} terdistribusi berbasis CRDT dengan
    \emph{handoff}. \emph{Contoh:} memindahkan \emph{process} ke node tersisa saat node hilang.
  \item[Handoff] Pemindahan \emph{process} beserta \emph{state}-nya ke node lain.
    \emph{Contoh:} Horde memindahkan \emph{process} dokumen sebelum \emph{pod} mati.
  \item[Phoenix] Kerangka kerja web Elixir. \emph{Contoh:} \emph{endpoint} HTTP pada \emph{port} 4000.
  \item[\texttt{:global}] \emph{Registry} nama global lintas-node bawaan Erlang.
    \emph{Contoh:} mendaftarkan \emph{singleton} agar unik di seluruh \emph{cluster}.
  \item[CRDT] \emph{Conflict-free Replicated Data Type}---konvergen tanpa koordinasi.
    \emph{Contoh:} Horde memakai \texttt{DeltaCrdt}.
  \item[Erlang distribution cookie] Rahasia bersama yang harus seragam antar-node.
    \emph{Contoh:} disuntikkan via \emph{Secret} Kubernetes.
  \item[Kuorum BEAM] Jumlah node minimum agar satu sisi \emph{cluster} dianggap otoritatif
    (mayoritas). \emph{Contoh:} \texttt{Horde.UniformQuorumDistribution} mematikan sisi
    dengan kurang dari setengah node.
\end{description}

\section{Kubernetes}
\begin{description}[style=nextline,leftmargin=1.4em]
  \item[Kubernetes] Orkestrator kontainer yang menjadwalkan, memantau, dan memulihkan
    beban kerja. \emph{Contoh:} me-\emph{restart pod} yang gagal \emph{probe}.
  \item[Control plane] Otak Kubernetes (API \emph{server}, \emph{scheduler}, \emph{controller}).
    \emph{Contoh:} menyimpan keadaan ke etcd.
  \item[etcd] Penyimpanan \emph{key-value} konsisten (Raft), \emph{backend control plane}.
    \emph{Contoh:} mendukung penguncian \emph{Lease}.
  \item[Pod] Unit penyebaran terkecil (satu/lebih kontainer). \emph{Contoh:} satu \emph{pod}
    memuat satu instans BEAM.
  \item[Node] Mesin (fisik/VM) anggota \emph{cluster}. \emph{Contoh:} \emph{pod}
    dijadwalkan ulang ke node sehat saat node-nya mati.
  \item[kubelet] Agen Kubernetes di tiap node yang menjalankan \emph{probe} dan mengelola
    siklus hidup \emph{pod}. \emph{Contoh:} kubelet mengirim \texttt{SIGTERM} saat \emph{pod} diterminasi.
  \item[Probe] Pemeriksaan kesehatan berkala oleh kubelet (\emph{liveness}/\emph{readiness}/\emph{startup}).
    \emph{Contoh:} \texttt{HTTP GET /healthz}.
  \item[Liveness probe] ``Apakah kontainer hidup?''; gagal $\rightarrow$ di-\emph{restart}.
    \emph{Contoh:} memeriksa \emph{port} 4000.
  \item[Readiness] Keadaan ``siap menerima lalu lintas''; \emph{readiness probe} mengatur
    keanggotaan \emph{endpoint Service}. \emph{Contoh:} \texttt{/readyz} mengembalikan 503
    sampai \emph{cluster} terbentuk.
  \item[Startup probe] \emph{Probe} untuk boot lambat sebelum \emph{liveness} aktif.
    \emph{Contoh:} menunggu migrasi basis data.
  \item[failureThreshold] Jumlah kegagalan \emph{probe} berturut-turut sebelum kontainer
    dianggap gagal. \emph{Contoh:} \texttt{failureThreshold: 1} membunuh \emph{pod} setelah
    satu kegagalan (positif palsu pada Masalah~1).
  \item[Readiness gate] Kondisi \emph{pod} kustom tambahan yang memengaruhi \emph{readiness}.
    \emph{Contoh:} siap hanya bila anggota \emph{cluster} BEAM \texttt{Up}.
  \item[SIGTERM] Sinyal terminasi sopan yang dikirim sebelum penghentian paksa.
    \emph{Contoh:} memicu \emph{drain} dan \emph{handoff}.
  \item[SIGKILL] Sinyal penghentian paksa setelah \emph{grace period} habis (tak dapat
    ditangkap). \emph{Contoh:} membunuh \emph{process} yang belum selesai \emph{handoff}.
  \item[Graceful shutdown] Mematikan tertib: berhenti menerima koneksi, \emph{drain},
    \emph{handoff}, lalu berhenti. \emph{Contoh:} menyelesaikan permintaan yang sedang berjalan.
  \item[terminationGracePeriodSeconds (grace period)] Tenggat antara \texttt{SIGTERM} dan
    \texttt{SIGKILL}. \emph{Contoh:} \texttt{30} detik.
  \item[preStop hook] Perintah/HTTP sebelum \texttt{SIGTERM}. \emph{Contoh:} tidur sejenak
    agar penghapusan \emph{endpoint} merambat.
  \item[Blocking termination] Menahan penghapusan \emph{pod} (mis.\ via \emph{finalizer})
    hingga syarat terpenuhi. \emph{Contoh:} menunda terminasi sampai \emph{handoff} Horde selesai.
  \item[Drain / eviction] Pengosongan \emph{pod} dari node (gangguan sukarela).
    \emph{Contoh:} \texttt{kubectl drain} saat pemeliharaan.
  \item[Rolling update] Pembaruan bertahap yang mengganti \emph{pod} sedikit demi sedikit.
    \emph{Contoh:} mengganti \emph{image} versi baru tanpa \emph{downtime}.
  \item[PodDisruptionBudget (PDB)] Batas \emph{pod} yang boleh tidak tersedia saat gangguan
    sukarela. \emph{Contoh:} \texttt{minAvailable: 2}.
  \item[maxUnavailable] Jumlah/persentase \emph{pod} maksimum yang boleh tidak tersedia saat
    \emph{rolling update}/PDB. \emph{Contoh:} \texttt{maxUnavailable: 1}.
  \item[anti-affinity] Aturan penjadwalan agar \emph{pod} tidak ditempatkan berdekatan
    (mis.\ node sama). \emph{Contoh:} menyebar replika antar-node demi ketahanan.
  \item[backoff] Penundaan yang meningkat antar-percobaan \emph{restart}
    (\texttt{CrashLoopBackOff}). \emph{Contoh:} jeda kian lama saat \emph{pod} terus gagal.
  \item[Operator / Controller / CRD] Pengendali kustom yang merekonsiliasi \emph{Custom
    Resource}. \emph{Contoh:} operator menyediakan basis data per-penyewa.
  \item[Bonny] Kerangka menulis \emph{operator} Kubernetes dalam Elixir. \emph{Contoh:}
    mengawasi CRD dan merekonsiliasi keadaan.
  \item[Kubernetes Lease] Objek koordinasi (\texttt{kind: Lease}) didukung etcd untuk
    penguncian/\emph{leader election}/\emph{fencing}. \emph{Contoh:} Akka memakainya sebagai
    arbiter \emph{split-brain}~\cite{akkaLease}.
  \item[KEDA / HPA] Penskala otomatis berbasis metrik/event (KEDA) dan sumber daya (HPA).
    \emph{Contoh:} KEDA menskala dari panjang antrian Oban.
\end{description}

\section{Konsep, Evaluasi, dan Platform Lain}
\begin{description}[style=nextline,leftmargin=1.4em]
  \item[Chaos Mesh] Perkakas \emph{chaos engineering} Kubernetes. \emph{Contoh:}
    \texttt{NetworkChaos} membuat partisi; \texttt{PodChaos} mematikan \emph{pod}.
  \item[LitmusChaos] Kerangka \emph{chaos engineering} Kubernetes berbasis eksperimen.
    \emph{Contoh:} eksperimen \texttt{pod-delete}.
  \item[Fault injection] Penyuntikan kesalahan terkendali untuk menguji ketahanan.
    \emph{Contoh:} mematikan node saat beban penuh.
  \item[MAPE-K] \emph{Loop} otonomik \emph{Monitor-Analyze-Plan-Execute} di atas
    \emph{Knowledge}. \emph{Contoh:} operator sebagai \emph{self-healing loop}.
  \item[Split-brain] Partisi jaringan yang membuat dua sisi \emph{cluster} sama-sama aktif.
    \emph{Contoh:} dua sisi mendaftarkan \emph{singleton} yang sama.
  \item[Fencing] Mencegah sisi yang tidak otoritatif terus bertindak.
    \emph{Contoh:} hanya pemegang \emph{Lease} yang boleh bertahan.
  \item[Arbiter] Pihak eksternal yang memutuskan sisi mana yang otoritatif saat
    konflik/partisi. \emph{Contoh:} \emph{Kubernetes Lease} (etcd) menjadi \emph{arbiter}
    \emph{split-brain} pada M5.
  \item[Double recovery] Dua lapis memulihkan satu kegagalan secara mubazir.
    \emph{Contoh:} OTP me-\emph{restart} \emph{process} sembari K8s me-\emph{restart pod}.
  \item[Drift (konfigurasi)] Divergensi pengaturan antar sumber/lapis yang seharusnya
    konsisten. \emph{Contoh:} \emph{probe} K8s dan \texttt{max\_restarts} OTP tak sinkron.
  \item[MTTR] \emph{Mean Time To Recovery}---rerata waktu pulih. \emph{Contoh:} metrik
    utama evaluasi (Bab~\ref{ch:eval}).
  \item[kind] \emph{Kubernetes IN Docker}---\emph{cluster} lokal untuk pengujian.
    \emph{Contoh:} \emph{testbed} evaluasi.
  \item[Akka / Orleans] \emph{Toolkit actor} JVM (Akka) dan kerangka \emph{virtual actor}
    .NET (Orleans). \emph{Contoh:} Akka memantulkan keanggotaan \emph{cluster} ke
    \emph{readiness}~\cite{akkaHealth}.
\end{description}

\section{Istilah pada Diagram (Sekuens/Aktivitas)}
Istilah-istilah berikut muncul pada label dan catatan di dalam diagram.
\begin{description}[style=nextline,leftmargin=1.4em]
  \item[Cowboy / Bandit] Server HTTP untuk BEAM yang dipakai Phoenix. \emph{Contoh:}
    Cowboy menerima koneksi TCP di \emph{port} 4000.
  \item[Finch] Klien HTTP berbasis \emph{pool} koneksi. \emph{Contoh:} \emph{pool} Finch
    memanggil API \emph{fraud} eksternal pada \emph{PayGate}.
  \item[Ecto] Pustaka basis data Elixir dengan \emph{pool} koneksi. \emph{Contoh:}
    \emph{pool} Ecto menyambung ulang saat \emph{failover} basis data.
  \item[MatchError] Galat Elixir saat pencocokan pola gagal. \emph{Contoh:} \emph{frame}
    cacat memicu \texttt{MatchError} pada \emph{decoder}.
  \item[\texttt{:one\_for\_one}] Strategi supervisi: hanya anak yang gagal yang disalakan
    ulang. \emph{Contoh:} lihat Listing~\ref{lst:sup}.
  \item[EXIT] Sinyal yang dikirim saat \emph{process} berhenti/\emph{crash}, ditangkap
    supervisor. \emph{Contoh:} \emph{decoder} mengirim \texttt{EXIT} ke supervisornya.
  \item[nodedown] Peristiwa saat sebuah node distribusi terputus. \emph{Contoh:} saat
    partisi, tiap sisi melihat sisi lain \texttt{nodedown}.
  \item[Wedged (macet)] \emph{Process} yang hidup tetapi tak membuat kemajuan (mis.\
    terblokir tanpa \emph{timeout}). \emph{Contoh:} pekerja \emph{pool} terblokir pada
    koneksi yang menggantung.
  \item[Wedged VM] Instans BEAM yang \emph{process} OS-nya hidup tetapi tak melayani
    permintaan karena subsistem kritis terblokir. \emph{Contoh:} \emph{probe} TCP lulus
    padahal 100\% pembayaran macet (negatif palsu pada M1/M9).
  \item[Naive TCP probe] \emph{Probe} yang hanya memeriksa apakah \emph{port} menerima
    koneksi TCP, tanpa menilai kesehatan aplikasi sebenarnya. \emph{Contoh:}
    \texttt{tcpSocket: 4000} lulus meski seluruh pekerja \emph{pool} macet.
  \item[Scheduler / utilisasi scheduler] Penjadwal BEAM yang menjalankan \emph{process};
    utilisasinya menandakan beban nyata. \emph{Contoh:} salah satu sinyal M9.
  \item[Run-queue] Antrian \emph{process} siap-jalan per \emph{scheduler}; yang membengkak
    menandakan kemacetan. \emph{Contoh:} indikator kesehatan pada M9.
  \item[Mailbox / message-queue backpressure] Antrian pesan tiap \emph{process}; backlog
    yang tumbuh menandakan \emph{backpressure}. \emph{Contoh:} \emph{mailbox} membengkak
    saat \emph{process} tak mengimbangi laju pesan.
  \item[Restart intensity] Laju \emph{restart} anak oleh supervisor; lonjakannya menandakan
    kegagalan berulang. \emph{Contoh:} dipakai M1 dan M9.
  \item[Deroute] Mengeluarkan \emph{pod} dari \emph{Service Endpoints} sehingga lalu lintas
    berhenti diarahkan ke sana (tanpa \emph{restart}). \emph{Contoh:} \emph{readiness} gagal
    $\rightarrow$ \emph{pod} di-\emph{deroute}.
  \item[Per-subtree restart intensity] Intensitas \emph{restart} dihitung per subpohon
    supervisi, bukan secara global. \emph{Contoh:} M1 memeriksa tiap subpohon.
  \item[Wedged-child detection] Deteksi anak supervisi yang hidup tetapi macet.
    \emph{Contoh:} dasar vonis tidak-sehat pada M1/M9.
  \item[Critical subtree] Subpohon supervisi yang menjalankan fungsi penting;
    kemacetan/over-\emph{restart}-nya menentukan vonis \emph{probe}. \emph{Contoh:} M1.
  \item[Up/down (status anak)] Keadaan anak supervisi hidup (\emph{up}) atau mati
    (\emph{down}). \emph{Contoh:} dimasukkan ke perhitungan kesehatan M1.
  \item[\texttt{/healthz} (derived)] \emph{Endpoint} \emph{healthz} yang vonisnya diturunkan
    dari keadaan \emph{supervision tree} (M1), bukan pemeriksaan dangkal, sehingga Kubernetes
    bertindak (\emph{deroute}/\emph{restart}) atas kesehatan supervisi yang sebenarnya.
  \item[Health deriver] Komponen M1 yang memeriksa \emph{OTP supervision tree} dan menghitung
    vonis \emph{probe}. \emph{Contoh:} memetakan intensitas \emph{restart}/kemacetan subpohon
    menjadi balasan 200 atau 503.
  \item[Fault rate] Laju kegagalan teramati per satuan waktu, menjadi masukan adaptasi M8.
    \emph{Contoh:} lonjakan \emph{crash} \emph{decoder} menaikkan \emph{fault rate}.
  \item[Adaptive budget] Penyetelan dinamis jatah \emph{restart}
    (\texttt{max\_restarts}/\texttt{max\_seconds}) menurut kondisi, bukan nilai statis.
    \emph{Contoh:} jatah diperketat saat \emph{fault rate} tinggi.
  \item[\texttt{f(rate, history)}] Fungsi adaptasi yang memetakan \emph{fault rate} dan
    riwayat ke jatah \emph{restart} baru. \emph{Contoh:} inti \textsc{AdaptBudget} pada
    Algoritma~\ref{alg:m8}.
  \item[Child crash] Peristiwa anak supervisor berhenti tak normal. \emph{Contoh:} memicu
    evaluasi apakah masih dalam jatah \emph{restart}.
  \item[Escalate (eskalasi)] Menyerahkan kegagalan ke lapis di atas
    (anak~$\rightarrow$~supervisor~$\rightarrow$~VM~$\rightarrow$~Kubernetes). \emph{Contoh:}
    jatah habis $\Rightarrow$ VM keluar $\Rightarrow$ K8s me-\emph{restart pod}.
  \item[Boundary (batas kepemilikan)] Ambang yang memisahkan ``OTP menangani'' dari ``K8s
    me-\emph{restart pod}''. \emph{Contoh:} ditetapkan secara adaptif oleh M8.
  \item[Principled \& load-adaptive] Sifat batas eskalasi: berlandas prinsip/terukur dan
    menyesuaikan beban, bukan konstanta sembarang. \emph{Contoh:} batas OTP$\leftrightarrow$K8s pada M8.
  \item[Signal aggregator] Komponen M9 yang merangkum sinyal runtime BEAM menjadi satu
    skor. \emph{Contoh:} menggabungkan \emph{scheduler}, \emph{run-queue}, dan \emph{mailbox}.
  \item[Sinyal liveness komposit BEAM] Sinyal kesehatan M9 yang menggabungkan beberapa
    metrik runtime BEAM (\emph{scheduler}, \emph{run-queue}, \emph{mailbox}, \emph{restart
    intensity}) menjadi satu vonis \emph{liveness}, menggantikan \emph{probe} TCP/HTTP.
    \emph{Contoh:} keluaran \emph{signal aggregator} pada \texttt{/alive}.
  \item[Composite progress score] Skor kesehatan gabungan dari beberapa sinyal; di bawah
    ambang berarti tidak sehat. \emph{Contoh:} keluaran \emph{signal aggregator} M9.
  \item[/alive (composite)] \emph{Endpoint} \emph{liveness} yang melayani skor komposit,
    bukan sekadar \emph{probe} TCP. \emph{Contoh:} dipakai M9 menggantikan \emph{probe} naif.
  \item[\texttt{terminate/2}, \texttt{handle\_continue/2}] \emph{Callback} OTP untuk
    menyimpan \emph{state} saat berhenti dan memulihkannya setelah inisialisasi.
    \emph{Contoh:} \emph{handoff} Horde menyimpan \emph{state} via \texttt{terminate/2}.
  \item[Phoenix.PubSub] \emph{Pub/sub} terdistribusi lintas-node. \emph{Contoh:} menyiarkan
    pesan ke semua pengguna antar-\emph{pod}.
  \item[Service / Endpoints] Abstraksi K8s yang merutekan lalu lintas ke \emph{pod} siap;
    daftar \emph{Endpoints} berisi \emph{pod} anggota. \emph{Contoh:} \emph{pod} ditambahkan
    ke \emph{Endpoints} saat \emph{ready}.
  \item[Deployment / ReplicaSet] Pengendali K8s yang menjaga jumlah replika dan melakukan
    \emph{rolling update}. \emph{Contoh:} Deployment mengganti \emph{pod} secara bertahap.
  \item[Voluntary disruption (gangguan sukarela)] Gangguan yang dipicu operator
    (\emph{drain}, \emph{rolling update}, \emph{scale-in})---bukan kegagalan---yang dibatasi
    PDB. \emph{Contoh:} \texttt{kubectl drain} saat pemeliharaan node.
  \item[Cluster change] Perubahan keanggotaan \emph{cluster} (node bergabung/keluar) yang
    memicu perhitungan ulang PDB/kuorum pada M7. \emph{Contoh:} node baru
    $\rightarrow$ hitung ulang \texttt{minAvailable}.
  \item[minAvailable] Parameter PDB: jumlah/persentase \emph{pod} minimum yang harus tetap
    tersedia. \emph{Contoh:} \texttt{minAvailable} disetel sama dengan kuorum BEAM.
  \item[Allow / block eviction] Keputusan API saat permintaan pengusiran \emph{pod}:
    \emph{allow} bila PDB terpenuhi, \emph{block} bila melanggar kuorum. \emph{Contoh:}
    logika M7.
  \item[Handoff throughput] Laju Horde memindahkan \emph{process}/\emph{state} per satuan
    waktu; membatasi berapa \emph{pod} boleh \emph{down} serentak. \emph{Contoh:}
    \texttt{maxUnavailable} diturunkan darinya.
  \item[Handoff backlog] Jumlah \emph{process}/\emph{state} yang masih menunggu di-\emph{handoff}.
    \emph{Contoh:} masukan M3 untuk menghitung \emph{grace}.
  \item[libcluster convergence] Waktu sampai libcluster menemukan seluruh \emph{peer} dan
    \emph{cluster} stabil. \emph{Contoh:} M3 menambahkannya ke perhitungan \emph{grace}.
  \item[In-flight work] Permintaan/tugas yang sedang diproses saat \emph{shutdown} dimulai.
    \emph{Contoh:} harus di-\emph{drain}/diselesaikan sebelum berhenti.
  \item[Required grace] \emph{Grace period} minimum agar \emph{handoff} dan \emph{drain}
    tuntas, dihitung M3. \emph{Contoh:} $\max(\textit{backlog}/\textit{throughput},\ \textit{convergence},\ g_{\min})$.
  \item[Rollout pace] Laju penggantian \emph{pod} saat \emph{rolling update}, diatur
    \texttt{maxUnavailable}. \emph{Contoh:} diperlambat M3 bila \emph{handoff} melambat.
  \item[One fault $\rightarrow$ one recovery authority] Prinsip M2: tiap kegagalan ditangani
    tepat oleh satu lapis (OTP \emph{atau} Kubernetes), mencegah \emph{double recovery}.
    \emph{Contoh:} kegagalan \emph{process} dimiliki OTP; matinya node dimiliki Kubernetes.
  \item[Failure classifier / ownership policy] Komponen M2 yang mengklasifikasikan kelas
    kegagalan lalu menerapkan kebijakan kepemilikan untuk memutuskan lapis pemulihan.
    \emph{Contoh:} mengklasifikasi \emph{transient}/\emph{resource}/\emph{node} dan
    mengarahkan ke OTP atau Kubernetes.
  \item[Fault event] Peristiwa kegagalan (\emph{crash}, \texttt{EXIT}, galat dependensi)
    yang memicu pengklasifikasi M2. \emph{Contoh:} sinyal \emph{crash} \emph{process} masuk
    ke \emph{failure classifier}.
  \item[\texttt{:global}/Horde quorum] Kuorum yang harus dipertahankan agar \texttt{:global}
    dan Horde tetap otoritatif. \emph{Contoh:} gangguan sukarela tak boleh memutusnya.
  \item[BEAM metrics] Metrik runtime BEAM (utilisasi \emph{scheduler}, \emph{run-queue},
    \emph{mailbox}, \emph{restart intensity}) yang diekspor untuk \emph{probe} atau
    penskalaan. \emph{Contoh:} sumber sinyal M9 dan KEDA.
  \item[CrashLoopBackOff] Status \emph{pod} K8s saat kontainer berulang gagal mulai; jeda
    \emph{backoff} meningkat tiap percobaan. \emph{Contoh:} eskalasi M2/M8 tampak sebagai
    CrashLoopBackOff.
  \item[Quorum controller] Komponen M7 yang menghitung kuorum dan kapasitas \emph{handoff}
    lalu menyetel PDB. \emph{Contoh:} menetapkan \texttt{minAvailable} = \emph{quorum threshold}.
  \item[Cluster size] Jumlah node anggota \emph{cluster} BEAM saat ini. \emph{Contoh:}
    masukan untuk menghitung kuorum.
  \item[Quorum threshold] Ambang minimum node agar kuorum terpenuhi (mis.\ mayoritas).
    \emph{Contoh:} $\lfloor n/2 \rfloor + 1$.
  \item[K8s API (API server)] Antarmuka \emph{control plane} Kubernetes untuk membaca/mengubah
    objek. \emph{Contoh:} \emph{operator} menulis PDB lewat K8s API.
  \item[Eviction / drain API] API Kubernetes untuk mengusir/mengosongkan \emph{pod} yang
    menghormati PDB. \emph{Contoh:} \emph{drain} memanggil \emph{Eviction API} yang dapat
    ditolak PDB.
  \item[Cross-layer policy (kebijakan lintas-lapis)] Kebijakan pemulihan yang konsisten di
    kedua lapis (supervisi OTP dan manifes Kubernetes), idealnya dari satu sumber kebenaran.
    \emph{Contoh:} M6 mengompilasi satu spesifikasi menjadi kebijakan yang konsisten
    \emph{by construction}.
  \item[K8s manifest (manifes K8s)] Berkas deklaratif (YAML) yang mendefinisikan objek
    Kubernetes (\emph{Deployment}, \emph{Service}, PDB, \emph{probe}). \emph{Contoh:} M6
    membangkitkan manifes lengkap dengan \emph{probe}, PDB, \emph{grace}, dan \emph{anti-affinity}.
  \item[Change the spec once (ubah spesifikasi sekali)] Prinsip M6: mengubah satu
    spesifikasi sumber memicu pembangkitan ulang kedua lapis sehingga mencegah \emph{drift}.
    \emph{Contoh:} ubah strategi pemulihan di satu tempat $\rightarrow$ spesifikasi OTP dan
    manifes K8s diperbarui serempak.
  \item[Single FT spec (spesifikasi FT tunggal)] Satu berkas spesifikasi \emph{fault
    tolerance} deklaratif (sifat \emph{stateful}, strategi pemulihan, kuorum, \emph{timing})
    yang menjadi sumber kebenaran tunggal pada M6; pengembang \emph{menulis}-nya sekali
    (\emph{author single FT spec}) lalu ia dikompilasi ke spesifikasi OTP dan manifes K8s.
    \emph{Contoh:} satu deklarasi menggantikan konfigurasi terpisah di kedua lapis.
  \item[ClusterMembershipCheck] Pemeriksaan \emph{readiness} Akka berbasis keanggotaan
    \emph{cluster}~\cite{akkaHealth}. \emph{Contoh:} \emph{pod} siap hanya bila anggotanya \texttt{Up}.
  \item[BEAM operator (Bonny)] \emph{Operator} Kubernetes yang ditulis dalam Elixir memakai
    Bonny dan sadar keanggotaan \emph{cluster} BEAM (Horde/distribusi Erlang). \emph{Contoh:}
    M4 memantulkan keanggotaan ke \emph{readiness} dan menahan terminasi \emph{pod}.
  \item[Supervision-awareness] Sifat M4: \emph{readiness}/\emph{probe} diturunkan dari keadaan
    \emph{supervision tree}, bukan sekadar keanggotaan datar. \emph{Contoh:} subpohon kritis
    yang macet membuat \emph{pod} ditandai tidak siap.
  \item[Handoff-blocking termination] Menahan terminasi \emph{pod} (via \emph{finalizer}/\emph{preStop})
    hingga \emph{handoff} Horde selesai. \emph{Contoh:} bagian M4 yang baru bagi BEAM.
  \item[Erlang distribution] Protokol bawaan yang menghubungkan node BEAM menjadi
    \emph{cluster} (koneksi node, pemantauan, \texttt{:global}, peristiwa \texttt{nodedown}).
    \emph{Contoh:} libcluster membentuknya; M4 mengawasi keanggotaannya.
  \item[lease-majority] Strategi \emph{Split Brain Resolver} Akka yang memakai
    \emph{Kubernetes Lease} sebagai arbiter~\cite{akkaSBR}. \emph{Contoh:} hanya sisi pemegang
    \emph{lease} yang bertahan.
  \item[compare-and-swap / etcd CAS] Operasi atomik pada etcd yang menjamin hanya satu
    penulis berhasil; mendasari pengambilan \emph{Kubernetes Lease} (\emph{K8s lease}).
    \emph{Contoh:} saat partisi, hanya satu sisi memenangkan \emph{etcd CAS} dan memegang
    \emph{lease}.
  \item[Keeps singletons (mempertahankan singleton)] Sisi otoritatif (pemegang \emph{lease})
    terus menjalankan \emph{singleton}-nya tanpa gangguan. \emph{Contoh:} partisi pemenang
    tetap memegang papan peringkat.
  \item[Fences itself (mem-fence diri)] Sisi yang kalah \emph{lease} menghentikan dirinya
    agar tak bertindak otoritatif. \emph{Contoh:} partisi minoritas berhenti menjalankan
    \emph{singleton}-nya.
  \item[Stops minority / no duplicate] Menghentikan sisi minoritas sehingga hanya satu
    instans \emph{singleton} yang aktif---mencegah duplikat dan \emph{last-writer-wins}.
    \emph{Contoh:} hasil \emph{fencing} pada M5.
  \item[PromEx / Prometheus] PromEx mengekspor metrik Elixir; Prometheus mengikis dan
    menyimpannya~\cite{keda}. \emph{Contoh:} metrik \emph{run-queue} diekspor untuk penskalaan.
\end{description}
